\begin{abstract}
We study the problem of sequential dynamic model ensembling and model merging over deep networks,
which has emerged as a key paradigm for multi-task model serving. Gating routing coefficients
in sequential architectures often suffer from violent, high-frequency oscillations across
network depth---a phenomenon we formalize as \emph{sequential routing jitter}---and show
severe susceptibility to transductive overfitting under extreme data scarcity. Prior efforts
address this heuristically via empirical trajectory smoothing, which lacks rigorous
stability guarantees. In this work, we propose the \textbf{Contraction-Regularized Router (CR-Router)},
formulating sequential deep ensembling as a discrete-time dynamical system. Under the framework
of Banach's Fixed-Point Theorem, we prove a novel Lipschitz bound on the joint layer-wise
representation-routing map. By introducing explicit spectral norm regularization and inverse
routing temperature penalties, we constrain the system mapping to be a strict contraction,
theoretically guaranteeing convergence to a unique fixed-point trajectory under depth.

Empirically, we evaluate CR-Router inside a 14-layer, 192-dimensional Analytical Coordinate
Sandbox across 10 random seeds under extreme data scarcity (16 calibration samples per task).
Under perfectly orthogonal task subspaces (Experiment 1), CR-Router achieves a stellar
classification accuracy of \textbf{53.35\% $\pm$ 3.84\%} and routing accuracy of
\textbf{65.10\% $\pm$ 2.70\%} (representing an \textbf{18.62\%} absolute improvement over the
unregularized router). Under overlapping task subspaces (Experiment 2), static Uniform Merging
collapses to 27.48\% due to representation cross-talk, whereas CR-Router recovers a strong
classification accuracy of \textbf{43.48\% $\pm$ 4.70\%} (outperforming static Uniform Merging
by \textbf{16.00\%} and unregularized routing by \textbf{12.86\%} in absolute performance), while
validating our theoretical stability predictions. Finally, we propose \emph{Adaptive Test-Time Temperature Annealing}
to decouple training-time optimization stability from inference-time sharpness, unlocking substantial additional gains.
\end{abstract}